% Base tcolorbox options shared by all theorem boxes.
% Uses only standard tcolorbox options (no enhanced/skins) for Beamer compatibility.
% leftrule=3pt + colframe (set per style) replaces enhanced+frame hidden+borderline west.
\tcbset{
  adenc thmboxbase/.style = {
    breakable,
    arc = 0pt,
    leftrule = 3pt,
    rightrule = 0pt,
    toprule = 0pt,
    bottomrule = 0pt,
    left = .75em,
    right = 0pt,
    top = 3pt,
    bottom = 5pt,
    before skip = 5pt,
    after skip = .5em,
  }
}

% Creates a thmtools style (head formatting only) and a named tcolorbox style.
% (usage: \getThmboxStyle{style name}{foreground color}{background color}{body font})
\newcommand{\getThmboxStyle}[4]{
  \declaretheoremstyle[
    headfont = \bfseries\color{#2},
    notefont = \bfseries\color{#2},
    bodyfont = #4,
    headpunct = {. },
    postheadspace = {0pt},
  ]{#1}
  \tcbset{adenc #1/.style = {adenc thmboxbase, colback=#3, colframe=#2}}
}

% Colors for theorem environments
\definecolor{color1}{HTML}{052E66} % blue, thmbox
\definecolor{color2}{HTML}{8F0A00} % red, defbox
\definecolor{color3}{HTML}{2B4E2C} % green, exbox
\definecolor{color5}{HTML}{764506} % orange, probbox

% If color passed, use colors; otherwise black and white
\adencIfColor{%
  \getThmboxStyle{thmbox}{color1}{color1!5}{\itshape}
  \getThmboxStyle{defbox}{color2}{color2!5}{\normalfont}
  \getThmboxStyle{exbox}{color3}{color3!5}{\normalfont}
  \getThmboxStyle{probbox}{color5}{color5!5}{\normalfont}
}{%
  \getThmboxStyle{thmbox}{black}{white}{\itshape}
  \getThmboxStyle{defbox}{black}{white}{\normalfont}
  \getThmboxStyle{exbox}{black}{white}{\normalfont}
  \getThmboxStyle{probbox}{black}{white}{\normalfont}
}

% If theorembox=plain|color: use tcolorbox-boxed styles (breakable across pages), no end symbols
% Default (theorembox=none): use default styles `plain`, `definition`, and `remark`, with end symbols
\adencIfBoxed{%
  \newcommand{\thmboxStyle}{thmbox}
  \newcommand{\defboxStyle}{defbox}
  \newcommand{\exboxStyle}{exbox}
  \newcommand{\probboxStyle}{probbox}
  % No end symbols for example and remark
  \newcommand\remqed{}
  \newcommand\exqed{}
}{%
  \newcommand{\thmboxStyle}{plain}
  \newcommand{\defboxStyle}{definition}
  \newcommand{\exboxStyle}{remark}
  \newcommand{\probboxStyle}{remark}
  % End symbols for example and remark
  \adencIfFancyqed{%
    \newcommand\remqed{\Coffeecup}
    \newcommand\exqed{\Faxmachine} % \Telefon
  }{%
    \newcommand\remqed{\(\diamond\)}
    \newcommand\exqed{\(\triangleleft\)}
  }
}

\newcommand\proofqedsymbol{\qedsymbol}


\declaretheorem[style = \thmboxStyle, name = Theorem, numberwithin = section]{theorem}
\declaretheorem[style = \thmboxStyle, name = Theorem, numbered = no]{theorem*}
\declaretheorem[style = \thmboxStyle, name = Lemma, sibling = theorem]{lemma}
\declaretheorem[style = \thmboxStyle, name = Lemma, numbered = no]{lemma*}
\declaretheorem[style = \thmboxStyle, name = Proposition, sibling = theorem]{proposition}
\declaretheorem[style = \thmboxStyle, name = Proposition, numbered = no]{proposition*}
\declaretheorem[style = \thmboxStyle, name = Corollary, sibling = theorem]{corollary}
\declaretheorem[style = \thmboxStyle, name = Corollary, numbered = no]{corollary*}

\declaretheorem[style = \defboxStyle, name = Definition, sibling = theorem]{definition}
\declaretheorem[style = \defboxStyle, name = Definition, numbered = no]{definition*}

\declaretheorem[style = \exboxStyle, name = Example, sibling = theorem, qed = \exqed]{example}
\declaretheorem[style = \exboxStyle, name = Example, numbered = no, qed = \exqed]{example*}


% or use \noteboxStyle
\declaretheorem[style = remark, name = Remark, sibling = theorem, qed = \remqed]{remark}
\declaretheorem[style = remark, name = Remark, numbered = no, qed = \remqed]{remark*}
\declaretheorem[style = remark, name = Note, sibling=theorem]{note}
\declaretheorem[style = remark, name = Note, numbered=no]{note*}
\declaretheorem[style = remark, name = Assumption, sibling=theorem]{assumption}
\declaretheorem[style = remark, name = Assumption, numbered=no]{assumption*}
\declaretheorem[style = remark, name = Intuition, sibling=theorem]{intuition}
\declaretheorem[style = remark, name = Intuition, numbered=no]{intuition*}
\declaretheorem[style = remark, name = Properties, sibling=theorem]{properties}
\declaretheorem[style = remark, name = Properties, numbered=no]{properties*}
\declaretheorem[style = remark, name = Fact, sibling = theorem]{fact}
\declaretheorem[style = remark, name = Fact, numbered = no]{fact*}
\declaretheorem[style = remark, name = Claim, sibling = theorem]{claim}
\declaretheorem[style = remark, name = Claim, numbered = no]{claim*}
\declaretheorem[style = remark, name = Notation, sibling = theorem]{notation}
\declaretheorem[style = remark, name = Notation, numbered = no]{notation*}
\declaretheorem[style = remark, name = Conjecture, sibling = theorem]{conjecture}
\declaretheorem[style = remark, name = Conjecture, numbered = no]{conjecture*}

\declaretheorem[style = \probboxStyle, name = Problem, sibling = theorem]{problem}
\declaretheorem[style = \probboxStyle, name = Problem, numbered = no]{problem*}
\declaretheorem[style = \probboxStyle, name = Exercise, sibling = theorem]{exercise}
\declaretheorem[style = \probboxStyle, name = Exercise, numbered = no]{exercise*}
\declaretheorem[style = \probboxStyle, name = Question, sibling = theorem]{question}
\declaretheorem[style = \probboxStyle, name = Question, numbered = no]{question*}

% for cref:
% \crefname{assumption}{Assumption}{Assumptions}
% \Crefname{assumption}{Assumption}{Assumptions}


% proof environment style:
\declaretheoremstyle[
  headfont=\bfseries,
  notefont = \normalfont\itshape,
  bodyfont=\normalfont,
  numbered=no,
  qed=\proofqedsymbol
]{thmproofbox}
\declaretheorem[numbered=no, style=thmproofbox, name=Proof]{replacementproof}

\makeatletter
\renewenvironment{proof}[1][\proofname]{%
    \ifx#1\proofname
        \begin{replacementproof}%
    \else
        \begin{replacementproof}[#1]%
    \fi
}{\end{replacementproof}}
\makeatother


% Apply tcolorbox wrapping for boxed mode.
% \tcolorboxenvironment wraps an existing environment in a tcolorbox using
% etoolbox hooks, preserving optional arguments like [note].
\adencIfBoxed{%
  \tcolorboxenvironment{theorem}{adenc thmbox}
  \tcolorboxenvironment{theorem*}{adenc thmbox}
  \tcolorboxenvironment{lemma}{adenc thmbox}
  \tcolorboxenvironment{lemma*}{adenc thmbox}
  \tcolorboxenvironment{proposition}{adenc thmbox}
  \tcolorboxenvironment{proposition*}{adenc thmbox}
  \tcolorboxenvironment{corollary}{adenc thmbox}
  \tcolorboxenvironment{corollary*}{adenc thmbox}
  \tcolorboxenvironment{definition}{adenc defbox}
  \tcolorboxenvironment{definition*}{adenc defbox}
  \tcolorboxenvironment{example}{adenc exbox}
  \tcolorboxenvironment{example*}{adenc exbox}
  \tcolorboxenvironment{problem}{adenc probbox}
  \tcolorboxenvironment{problem*}{adenc probbox}
  \tcolorboxenvironment{exercise}{adenc probbox}
  \tcolorboxenvironment{exercise*}{adenc probbox}
  \tcolorboxenvironment{question}{adenc probbox}
  \tcolorboxenvironment{question*}{adenc probbox}
}{}
